% GENERATED FILE. Edit canonical lessons, then run npm run generate:books.
% canonical-source-hash: fnv1a64:018d0826b712077b
% canonical-lessons: ZH-C81-li, ZH-W81-li-inside, ZH-C81-chengshi, ZH-C81-difang, ZH-C81-waiguo, ZH-C81-pangbian

\chapter{The City}
\label{ch:zh-the-city}
\hlchaptermodality{81}

\begin{chapteropening}
\textbf{By the end of this chapter:} \emph{I can talk about cities and places, say inside and beside, and talk about abroad.}

Complete the last lesson of chapter 81: I can talk about cities and places, say inside and beside, and talk about abroad.
\end{chapteropening}

\section[\texorpdfstring{\zh{里} (lǐ)}{里 (lǐ)}]{\zh{里} (lǐ) — inside; in}

\label{lesson:ZH-C81-li}



\begin{quote}
Before the new one: say the Mandarin for doctor, then the Mandarin for company.
\end{quote}

\subsection*{You'll want to know: \zh{里}}
\textbf{\zh{里}} — \emph{lǐ} — \textquotedblleft{}inside; in\textquotedblright{}.

\begin{quote}
\textbf{\zh{山里有马}} \emph{shān lǐ yǒu mǎ} — There are horses up in the hills.
\end{quote}

The new character, in pieces: \zh{里} — \zh{日} with a vertical running down through it, standing on two bars. The next lesson puts it in your hand.

\begin{cousinweb}
\textbf{\zh{里}} (\emph{lǐ}) is inherited from Middle Chinese \emph{liX} (Guangyun, in Baxter's transcription), glossed 'inside' and, through an older book, 'the inside of a garment'.
\end{cousinweb}

\begin{grammarlens}[title={What you've built}]
One more word for meeting people.
\end{grammarlens}

\subsection*{Your turn}
\begin{itemize}
\raggedright
  \item \emph{Say it:} \emph{lǐ}
  \item \emph{Say it:} \emph{lǐ}, once more
  \item \emph{Say it:} \emph{shān lǐ yǒu mǎ} — There are horses up in the hills.
  \item \emph{Recall:} say the Mandarin for doctor, then the Mandarin for company, then say \emph{lǐ} again
\end{itemize}

\begin{tcolorbox}[breakable,colback=teal!4,colframe=teal!35!black,title={Before you move on}]
What does \zh{里} mean? (\textquotedblleft{}Inside; in\textquotedblright{}.) Say it once more.
\end{tcolorbox}
\section[\texorpdfstring{\zh{里} (lǐ)}{里 (lǐ)}]{\zh{里} — \zh{日}, a line through it, two bars}

\label{lesson:ZH-W81-li-inside}



\begin{quote}
The word first, then the character inside it.

\begin{itemize}
\raggedright
  \item \emph{Recall:} say \emph{lǐ} — \textbf{R1}, one lesson back, before it decays
\end{itemize}

\begin{quote}
\zh{日} \zh{家}
\end{quote}

\emph{Read it:} each one aloud, with nothing said first — \emph{rì}, \emph{jiā}
\end{quote}

\subsection*{Script — \zh{日}, a line through it, two bars}
\hlblockfigure{\detokenize{figures/ZH-W81-li-inside-filmstrip.pdf}}{How \zh{里} is written, stroke by stroke}

\begin{quote}
\zh{里}
\end{quote}

\textbf{One piece, built from the top down.}

\textbf{\zh{日}, four strokes}: the left side, the top that turns down, the middle bar, the bottom bar.

\textbf{A vertical} down the middle, through \zh{日} and on below it.

\textbf{Two horizontals}, a short one across the vertical, then a long one along the bottom.

4 + 1 + 2 = seven.

\subsection*{Your turn}
\begin{itemize}
\raggedright
  \item \emph{Write it:} \zh{里}, all seven, with the model in front of you
  \item \emph{Cover:} write it again from memory, and count the pieces, not the strokes
\end{itemize}

\begin{tcolorbox}[breakable,colback=teal!4,colframe=teal!35!black,title={Before you move on}]
How many strokes in \zh{里}? (\textbf{Seven}.) Which character you know is at the top? (\textbf{\zh{日}}.) Which bar is longest? (\textbf{The last, along the bottom}.)

Source: \href{https://raw.githubusercontent.com/chanind/hanzi-writer-data/68d10a4b21150cae5e1ebbd223eed289cf32d90c/data/\%E9\%87\%8C.json}{Hanzi Writer Data \zh{里}, pinned stroke paths and medians}.
\end{tcolorbox}
\section[chéngshì]{chéngshì — city}

\label{lesson:ZH-C81-chengshi}



\begin{quote}
Before the new one: say the Mandarin for company, then the Mandarin for inside.
\end{quote}

\subsection*{You'll want to know: chéngshì}
\emph{chéngshì} — \textquotedblleft{}city\textquotedblright{}.

\begin{quote}
\emph{Běijīng shì yí ge dà chéngshì} — Beijing is a big city.
\end{quote}

Its characters come later; for now it is a word for the ear.

\begin{cousinweb}
\emph{chéngshì} is built from chéng 'city walls' and shì 'market', both inherited from Middle Chinese, \emph{dzyeng} and \emph{dzyiX} (Guangyun, in Baxter's transcription).
\end{cousinweb}

\begin{grammarlens}[title={What you've built}]
One more word for meeting people.
\end{grammarlens}

\subsection*{Your turn}
\begin{itemize}
\raggedright
  \item \emph{Say it:} \emph{chéngshì}
  \item \emph{Say it:} \emph{chéngshì}, once more
  \item \emph{Say it:} \emph{Běijīng shì yí ge dà chéngshì} — Beijing is a big city.
  \item \emph{Recall:} say the Mandarin for company, then the Mandarin for inside, then say \emph{chéngshì} again
\end{itemize}

\begin{tcolorbox}[breakable,colback=teal!4,colframe=teal!35!black,title={Before you move on}]
What does chéngshì mean? (\textquotedblleft{}City\textquotedblright{}.) Say it once more.
\end{tcolorbox}
\section[dìfang]{dìfang — place}

\label{lesson:ZH-C81-difang}



\begin{quote}
Before the new one: say the Mandarin for inside, then the Mandarin for city.
\end{quote}

\subsection*{You'll want to know: dìfang}
\emph{dìfang} — \textquotedblleft{}place\textquotedblright{}.

\begin{quote}
\emph{zhège dìfang hěn hǎo} — This is a nice place.
\end{quote}

Its characters come later; for now it is a word for the ear.

\begin{cousinweb}
\emph{dìfang} is built from dì 'land' and fāng 'quarter, region', both inherited from Middle Chinese, \emph{dijH} and \emph{pjang} (Guangyun, in Baxter's transcription).
\end{cousinweb}

\begin{grammarlens}[title={What you've built}]
One more word for meeting people.
\end{grammarlens}

\subsection*{Your turn}
\begin{itemize}
\raggedright
  \item \emph{Say it:} \emph{dìfang}
  \item \emph{Say it:} \emph{dìfang}, once more
  \item \emph{Say it:} \emph{zhège dìfang hěn hǎo} — This is a nice place.
  \item \emph{Recall:} say the Mandarin for inside, then the Mandarin for city, then say \emph{dìfang} again
\end{itemize}

\begin{tcolorbox}[breakable,colback=teal!4,colframe=teal!35!black,title={Before you move on}]
What does dìfang mean? (\textquotedblleft{}Place\textquotedblright{}.) Say it once more.
\end{tcolorbox}
\section[wàiguó]{wàiguó — foreign country; abroad}

\label{lesson:ZH-C81-waiguo}



\begin{quote}
Before the new one: say the Mandarin for city, then the Mandarin for place.
\end{quote}

\subsection*{You'll want to know: wàiguó}
\emph{wàiguó} — \textquotedblleft{}foreign country; abroad\textquotedblright{}.

\begin{quote}
\emph{tā shì wàiguó rén} — He's a foreigner.
\end{quote}

Its characters come later; for now it is a word for the ear.

\begin{cousinweb}
\emph{wàiguó} is built from wài 'outside' and \zh{国} 'country', both inherited from Middle Chinese, \emph{ngwajH} and \emph{kwok} (Guangyun, in Baxter's transcription).
\end{cousinweb}

\begin{grammarlens}[title={What you've built}]
One more word for meeting people.
\end{grammarlens}

\subsection*{Your turn}
\begin{itemize}
\raggedright
  \item \emph{Say it:} \emph{wàiguó}
  \item \emph{Say it:} \emph{wàiguó}, once more
  \item \emph{Say it:} \emph{tā shì wàiguó rén} — He's a foreigner.
  \item \emph{Recall:} say the Mandarin for city, then the Mandarin for place, then say \emph{wàiguó} again
\end{itemize}

\begin{tcolorbox}[breakable,colback=teal!4,colframe=teal!35!black,title={Before you move on}]
What does wàiguó mean? (\textquotedblleft{}Foreign country; abroad\textquotedblright{}.) Say it once more.
\end{tcolorbox}
\section[pángbiān]{pángbiān — side; beside}

\label{lesson:ZH-C81-pangbian}



\begin{quote}
Before the new one: say the Mandarin for place, then the Mandarin for abroad.
\end{quote}

\subsection*{You'll want to know: pángbiān}
\emph{pángbiān} — \textquotedblleft{}side; beside\textquotedblright{}.

\begin{quote}
\emph{wǒ jiā pángbiān yǒu yí ge shāngdiàn} — There's a shop next to my home.
\end{quote}

Its characters come later; for now it is a word for the ear.

\begin{cousinweb}
\emph{pángbiān} is built from páng 'side' and biān 'edge', both inherited from Middle Chinese, \emph{bang} and \emph{pen} (Guangyun, in Baxter's transcription).
\end{cousinweb}

\begin{grammarlens}[title={What you've built}]
One more word for meeting people.
\end{grammarlens}

\subsection*{Your turn}
\begin{itemize}
\raggedright
  \item \emph{Say it:} \emph{pángbiān}
  \item \emph{Say it:} \emph{pángbiān}, once more
  \item \emph{Say it:} \emph{wǒ jiā pángbiān yǒu yí ge shāngdiàn} — There's a shop next to my home.
  \item \emph{Recall:} say the Mandarin for place, then the Mandarin for abroad, then say \emph{pángbiān} again
\end{itemize}

\begin{tcolorbox}[breakable,colback=teal!4,colframe=teal!35!black,title={Before you move on}]
What does pángbiān mean? (\textquotedblleft{}Side; beside\textquotedblright{}.) Say it once more.
\end{tcolorbox}
